\textbf{T3: Dipole in a magnetic field (10 pts)}

Two small balls of mass $m$ each with charges $+q$ and $-q$ respectively,
connected by a rigid massless rod of length $d$, form a dipole. The dipole is
parallel to plane $XY$ and is placed in a uniform magnetic field $\vec{B}$
perpendicular to $XY$.

\begin{center}\includegraphics[width=0.32\linewidth]{figures/EuPhO_2022_Q3_fig1.png}\end{center}

Initially, the dipole is aligned with the direction $X$ and has initial angular
velocity $\omega_0$ in plane $XY$, as shown. Its center of mass is initially
located at origin and given initial velocity $\vec{v}_0$ parallel to $XY$, as
well.

Consider three distinct scenarios (a, b, c-d):

\textbf{(a)} \textit{(2 pts)} Find $\omega_0$ and the direction of $\vec{v}_0$,
so that the center of mass will move with the constant velocity
$\vec{v} = \vec{v}_0$?

\textbf{(b)} \textit{(3 pts)} Given $\omega_0$, find such $\vec{v}_0$
(direction and magnitude), so that the center of mass will travel in a circle.
Find the circle radius $R_c$ and the coordinates $x_c$, $y_c$ of its center.
You don’t need to prove the uniqueness of the solution.

\textbf{(c)} \textit{(4 pts)} Given $\vec{v}_0 = 0$, find the minimal
$\omega_0 = \omega_\mathrm{min}$ necessary for the dipole to reverse its
orientation during the motion.

\textbf{(d)} \textit{(1 pt)} If the dipole starts with $\vec{v}_0 = 0$ and
$\omega_0 = \omega_\mathrm{min}$ found in part (c), the trajectory of its
center of mass has an asymptote. Find the distance $D$ from the origin to the
asymptote.

\textit{Useful vector identity:}
\[
\vec{a} \times (\vec{b} \times \vec{c}) = \vec{b}\,(\vec{a} \cdot \vec{c}) - \vec{c}\,(\vec{a} \cdot \vec{b}),
\]
where “$\times$” and “$\cdot$” denote vector product and scalar product
respectively.
